\section{Version Control with Git}\label{sec:git}

The project's history is kept in a Git repository~\cite{progit} and published on GitHub at
\RepoURL, branch \texttt{main}. All commands were run in a terminal in the project folder.

\subsection{Setting up the repository}

The repository was created as soon as the project skeleton existed: the Maven wrapper, \texttt{pom.xml}, the
application class and the IDE's launch configuration. The steps were:
\begin{enumerate}
  \item \texttt{git init -b main}: create an empty repository whose first branch is \texttt{main};
  \item \texttt{git status}: check which files Git sees, all of them untracked at this point;
  \item \texttt{git add .}: stage the whole skeleton for the first commit, the build output and the local database
        being kept out by \texttt{.gitignore};
  \item \texttt{git status}: check that exactly the intended files are staged;
  \item \texttt{git commit}: record the initial commit;
  \item \texttt{git remote add origin} with the GitHub address, then \texttt{git remote -v} to check it;
  \item \texttt{git push -u origin main}: publish the branch and make \texttt{origin/main} its upstream.
\end{enumerate}
Listing~\ref{lst:01-initial-setup.txt} is the transcript of the first four steps as they ran.

\evidence[firstline=1, lastline=52]{01-initial-setup.txt}
  {Creating the repository and staging the skeleton}

The warnings printed by \texttt{git add} are Git for Windows noting that the files, written with Unix line endings,
will be checked out with Windows ones; they do not affect what is committed.

The last three steps are seen best in the repository as it is now (Listing~\ref{lst:remote}): the first commit
of the history is still the initial one, \texttt{origin} points at the GitHub repository, and \texttt{main} tracks
\texttt{origin/main}, with nothing committed that has not been pushed. (The untracked \texttt{report/} folder is
this report, before its own commit.)

\begin{lstlisting}[style=terminal, caption={The initial commit, the remote and the tracking branch, in the live
repository}, label={lst:remote}]
$ git log --reverse --format="%h %ad %s" --date=short | head -n 1
a0bb38c 2026-11-22 Initial commit: Spring Boot and JavaFX project skeleton on Maven

$ git remote -v
origin	https://github.com/hamdotexe/taskmanager.git (fetch)
origin	https://github.com/hamdotexe/taskmanager.git (push)

$ git status -sb
## main...origin/main
?? report/
\end{lstlisting}

\subsection{The workflow for each change}

Each piece of work, such as a new rule in the service, a screen, or a fix, was committed on its own, in this order:
\begin{enumerate}
  \item \texttt{git status}: see which files the change touched;
  \item \texttt{git diff}: read every changed line before staging it;
  \item \texttt{./mvnw -q test}: run the tests; nothing is committed unless they pass;
  \item \texttt{git add} followed by the files of the change, named one by one: staging by explicit path keeps
        unrelated edits and generated files out of the commit;
  \item \texttt{git commit -m "Area: what changed and why"}, with a second \texttt{-m} for a body when the reason
        needs more than a line;
  \item \texttt{git push}: publish the commit to \texttt{origin/main}.
\end{enumerate}
The history as a whole is reviewed with \texttt{git log --oneline --graph}, one line per commit.

\subsection{Commit messages}

Every message starts with the area of the program it changes, followed by what changed and, where it is not obvious,
why: for example ``Tasks: change a task's status on its own, as the list will'' or ``Test: tasks survive a restart of
the app''. A message that names its area makes the one-line log readable as a list of the program's parts as they
were built. When the reason needs more than a line, a body after a blank line gives it, as in this commit from the
history:
\begin{lstlisting}[style=terminal]
Tasks: trim the title and description, and refuse either over its limit

A title over 100 characters or a description over 1000 got past the service
and failed Bean Validation at commit, so the screen would have had a
TransactionSystemException to show. The service now refuses them with the
limit and the length typed, counted after trimming, so spaces around a text
never push it over. No description is kept as empty.

TaskRulesTest covers blank and missing titles, trimming, both limits at and
past the edge, a refused update leaving the task as it was, and a change
with no task chosen.
\end{lstlisting}

\subsection{Recording revisions: status, diff and log}

Three times during development, for a new rule and for two bug fixes, the output of \texttt{git status} and
\texttt{git diff} was saved with its commands just before the change was staged. In each, \texttt{git status} also
lists the \texttt{report/} folder as untracked: this report was being written at the same time and was committed on
its own.

\paragraph*{A new rule} A title over 100 characters or a description over 1000 used to get past the service and fail
Bean Validation at the commit, as described in Section~\ref{sec:design}. Listing~\ref{lst:02-status-diff-title-limits.txt}
shows the change that moved both limits into the service: \texttt{git status} lists the modified
\texttt{TaskService.java} and the new, still untracked test class \texttt{TaskRulesTest.java}, and \texttt{git diff}
shows the method \texttt{set} now trimming both texts and refusing either one over its limit, with the length typed in
the message.

\evidence[linerange={1-22,68-112}]{02-status-diff-title-limits.txt}
  {\texttt{git status} and \texttt{git diff} before the length limits were committed (the hunks that change only
   comments, lines 23--67 of the transcript, are left out)}

\paragraph*{A bug fix} With two copies of the program open on one file, a copy that had read a task could not save it
once the other copy had written in between: SQLite refused the write at once with \texttt{SQLITE\_BUSY\_SNAPSHOT},
the busy timeout did not help, and the save failed with a raw \texttt{CannotAcquireLockException} rather than a
message.
A new test, \texttt{TwoCopiesTest}, reproduced it with a second Spring context on the same database playing the other
copy. The fix is one setting in the connection URL (Listing~\ref{lst:03-status-diff-two-copies-fix.txt}): every
transaction now takes the write lock as it begins, so the other copy's change waits for this one to finish and then
goes through.

\evidence{03-status-diff-two-copies-fix.txt}
  {\texttt{git status} and \texttt{git diff} before the fix for two copies was committed}

\paragraph*{A bug fix in the window} When a dialog opened, the keyboard did not go to its first field: Ctrl+N
followed by typing lost the first keys, and the focus stayed in the list. The dimmed layer the dialog sits on is a
scroll pane, which puts its content in the window only once its skin is made, on the next pass of the stylesheet, so
asking the title field for the focus before then did nothing. The fix applies the stylesheet first
(Listing~\ref{lst:04-status-diff-dialog-focus-fix.txt}). This time \texttt{git diff} was given the one file of the
fix; \texttt{git status} shows the other files then changed, which went into commits of their own.

\evidence{04-status-diff-dialog-focus-fix.txt}
  {\texttt{git status} and \texttt{git diff} before the fix for the dialog's focus was committed}

Figures~\ref{fig:git-02-add-commit-push} to~\ref{fig:git-04-github} show the workflow in the terminal and on GitHub.
In Fig.~\ref{fig:git-02-add-commit-push}, a change to the README (the group's members) and to \texttt{pom.xml} (the
Spring Boot patch version, 4.1.2 to 4.1.3) is checked with \texttt{git status}, staged with \texttt{git add .}, which
with only those two files changed stages exactly them, committed and pushed. Fig.~\ref{fig:git-03-log} shows the
whole history as a graph, and Fig.~\ref{fig:git-04-github} the commits as GitHub lists them. The diffs themselves are
in the listings above.

\screenshotwide{git-02-add-commit-push}
  {A change checked with \texttt{git status} (README.md and pom.xml modified), then staged, committed and pushed
   with \texttt{git add}, \texttt{git commit} and \texttt{git push} (c1479d6..ae85057).}
  {Terminal: git add, git commit, git push of that change}

\screenshotwide{git-03-log}
  {The whole history with \texttt{git log --oneline --graph}: 22 commits, \texttt{HEAD} at ae85057.}
  {Terminal: \texttt{git log --oneline --graph} of the whole history}

\screenshotwide{git-04-github}
  {The repository's commits on GitHub.}
  {The repository's commits page on GitHub}

\subsection{The history}

When this report was built, the history had \CommitCount\ commits, made between \FirstCommitDate\ and
\LastCommitDate.
Table~\ref{tab:history} lists them oldest first, with the first line of each message, as
\texttt{git log --reverse} prints them.

\begin{table*}[!tp]
  \caption{The commit history}\label{tab:history}
  \centering\footnotesize
  \begin{tabularx}{\textwidth}{@{}rlL{0.15\textwidth}Y@{}}
    \toprule
    No. & Commit & Date & Message (first line) \\
    \midrule
    \CommitRows
    \bottomrule
  \end{tabularx}
\end{table*}
